\begin{afterword}
\summaryhead

The post is a thought experiment. The author offers three claims that sound crazy: a color of paint that reverses gravity, giant black spheres in the sky that lower male prostitutes for money, and grandchildren who think it evil to jail thieves instead of spanking them. Then the reader is placed in 1901 and asked to choose between those claims and three others: that there is an absolute speed limit at which length and time ``change around''; that a global network of adding machines will carry moving pictures of lesbian sex; and that grandchildren will think it evil to bar a black woman from the presidency. The second list is true. The idea comes from a comment by Robin Hanson, who wondered whether a detailed fiction could show how surprising the truth turned out to be.

\responsehead

The post makes one point and makes it briefly: some things that are true now would have sounded as crazy in 1901 as the post's inventions sound today. The three true claims support that. It is a companion to ``Making History Available,'' posted the day before, which asked readers to remember how often the world changed in ways no one guessed.

The comparison is built to come out one way. The true claims were chosen because they are strange, and they are worded in ways that make them sound stranger (``pretending they are made out of numbers''). The false claims were made up to be absurd. Such a pairing shows that the truth is sometimes strange, which is all the post claims. The post does not claim to show that, and it states no lesson at all.

The picture of 1901 is also simpler than the record. For physicists, the contraction of moving bodies was already in the literature, proposed by FitzGerald and by Lorentz to explain the Michelson--Morley result. And at least one popular forecast of 1900, by John Elfreth Watkins, described cameras connected electrically to screens thousands of miles away, and opened by saying its prophecies ``will seem strange, almost impossible.'' At least one writer of 1900 could picture part of this future, and expected it to sound strange.

In short: a short and fair illustration that the truth can sound crazy, built from lists chosen to make it so.
\end{afterword}
